\section{Round-11 Objective}
We pursue \textbf{(C) obstruction/correction}.  A concrete logical error entered the Round~7 and Round~8 proofs: they used the global inclusion
\[
A\subseteq A^{(k)},
\]
for fixed $k$, but this inclusion is false in general because $A$ does \emph{not} impose the $i$-th congruence condition below the threshold $n_i$, whereas $A^{(k)}$ does.

The new goal of this round is threefold.
\begin{enumerate}
\item Identify and prove the precise error.
\item Replace the false global inclusion by the correct \emph{tail inclusion}
\[
A\cap[n_k,\infty)\subseteq A^{(k)},
\]
and use it to derive a new unconditional theorem:
\[
\overline d(A)\le \delta,
\qquad
\overline\delta_{\log}(A)\le \delta,
\qquad
\delta:=\lim_{k\to\infty}d(A^{(k)}).
\]
This repairs the Round~8 zero-limit branch with a correct proof.
\item Combine this upper-density theorem with the Round~9/10 padding machinery to sharpen the obstruction: if a counterexample exists, then for every $\varepsilon>0$ there is one whose \emph{stage densities and actual upper/lower logarithmic densities all lie in a window of width $\varepsilon$ near $1$}.  Thus even \emph{small-amplitude oscillation near $1$} already contains the full difficulty of the problem.
\end{enumerate}

\section{Round-10 Foundation Used}
We use the following vetted earlier results.
\begin{itemize}
\item \textbf{Stage sets and stage densities.}  For each $k\ge 1$,
\[
A^{(k)}:=\{n\in\mathbb N:\ n\not\equiv a_i\pmod{n_i}\ \forall\ 1\le i\le k\}
\]
is periodic, so its natural density
\[
\delta_k:=d(A^{(k)})
\]
exists.

\item \textbf{Round~2 monotonicity.}  The stage sets decrease:
\[
A^{(k+1)}\subseteq A^{(k)},
\qquad
1\ge \delta_1\ge \delta_2\ge \cdots \ge 0,
\]
so $\delta_k\to \delta\in[0,1]$.

\item \textbf{Periodic asymptotics.}  If $P\subseteq\mathbb N$ is periodic with natural density $d(P)$, then
\[
\#\{n\le x:\ n\in P\}=d(P)x+O(1),
\qquad
\sum_{\substack{n\le x\\ n\in P}}\frac1n=d(P)\log x+O(1).
\]

\item \textbf{Round~5 single-progression harmonic estimate.}  For any modulus $m\ge 1$, any residue class $b\pmod m$, and integers $2\le X<Y$,
\[
\sum_{\substack{X\le n<Y\\ n\equiv b\ (\mathrm{mod}\ m)}}\frac1n
=
\frac1m\log\frac{Y}{X}+O\!\left(\frac1X\right),
\]
with an absolute implied constant.

\item \textbf{Round~9/10 affine padding formulas.}  If $(m_j,b_j)$ is an arbitrary instance with associated set $B$, and for $M\ge 3$ one forms the padded instance
\[
 n_1:=M,\ a_1:=0,
 \qquad
 n_{j+1}:=Mm_j,\ a_{j+1}:=Mb_j+1,
\]
then with $A_M$ the new set, one has exact affine relations
\begin{align*}
\underline{\delta}_{\log}(A_M)
&=
\frac{M-2}{M}+\frac1M\underline{\delta}_{\log}(B),
\\
\overline{\delta}_{\log}(A_M)
&=
\frac{M-2}{M}+\frac1M\overline{\delta}_{\log}(B),
\\
d\bigl(A_M^{(k+1)}\bigr)
&=
\frac{M-2}{M}+\frac1M d\bigl(B^{(k)}\bigr).
\end{align*}
In particular, logarithmic nonconvergence transfers exactly under padding.
\end{itemize}

\section{New Insight / Tool (Round-11)}
The correct replacement for the false global comparison $A\subseteq A^{(k)}$ is the \emph{tail comparison}
\[
A\cap[n_k,\infty)\subseteq A^{(k)}.
\]
This is enough to prove a new unconditional theorem:
\begin{equation}
\overline d(A)\le \delta,
\qquad
\overline\delta_{\log}(A)\le \delta,
\qquad
\delta:=\lim_{k\to\infty} d(A^{(k)}).
\tag{11.1}
\end{equation}
Thus the stage-limit $\delta$ is always an \emph{upper bound} for both the upper natural density and the upper logarithmic density of the true set $A$.

This immediately repairs the Round~8 zero-limit theorem: if $\delta=0$, then both upper densities vanish, hence $A$ has natural density $0$ and logarithmic density $0$.

Combining \eqref{11.1} with the Round~9/10 padding machinery yields a sharper obstruction than before: if a counterexample exists, then for every $\varepsilon>0$ there is one for which
\[
1-\varepsilon<\underline\delta_{\log}(A)<\overline\delta_{\log}(A)\le \delta<1,
\]
while simultaneously all stage densities stay within $\varepsilon$ of $\delta$ and the oscillation gap
\(\overline\delta_{\log}(A)-\underline\delta_{\log}(A)\)
is itself $<\varepsilon$.
So the full difficulty survives inside an \emph{arbitrarily thin oscillatory window near $1$}.

\section{Attack Plan (Round-11)}
The exact gap left after Round~10 was that the padding reduction showed the problem survives in an almost-stationary high-density branch, but it still used a Round~7/8 comparison whose proof was not valid.

We proceed as follows.
\begin{enumerate}
\item Exhibit a concrete counterexample to the false inclusion $A\subseteq A^{(k)}$.
\item Prove the corrected tail inclusion $A\cap[n_k,\infty)\subseteq A^{(k)}$.
\item Use periodic asymptotics of $A^{(k)}$ to prove the unconditional upper-density theorem \eqref{11.1}.
\item Recover the corrected zero-stage-density branch as an immediate corollary.
\item Repair the Round~7 tail criterion (hence also the summable-reciprocals theorem) using the corrected tail comparison.
\item Combine the new upper-density theorem with the Round~9/10 padding formulas to sharpen the obstruction: a counterexample, if it exists, can be taken to have tiny oscillation amplitude near $1$.
\end{enumerate}

\section{Work (Round-11)}
\subsection*{11.1. The false global inclusion and the correct tail inclusion}
\paragraph{Lemma 11.1 (the global inclusion used in Round~7/8 is false).}
The implication
\[
A\subseteq A^{(k)}
\]
is false in general, even for $k=1$.

\emph{Proof.}
Take $n_1=2$ and $a_1\equiv 1\pmod 2$.  Then $1\in A$ because $1<n_1$, so the first congruence condition is not active at $n=1$.  But $1\notin A^{(1)}$ because
\[
1\equiv 1\pmod 2.
\]
Hence $A\not\subseteq A^{(1)}$. \hfill $\square$

\paragraph{Lemma 11.2 (correct tail inclusion).}
For every fixed $k\ge 1$ and every integer $n\ge n_k$,
\begin{equation}
\mathbf 1_A(n)\le \mathbf 1_{A^{(k)}}(n).
\tag{11.2}
\end{equation}
Equivalently,
\begin{equation}
A\cap[n_k,\infty)\subseteq A^{(k)}.
\tag{11.3}
\end{equation}

\emph{Proof.}
If $n\ge n_k$ and $n\in A$, then for every $1\le i\le k$ we have $n_i\le n$, so the defining property of $A$ forces
\[
n\not\equiv a_i\pmod{n_i}
\qquad (1\le i\le k).
\]
That is exactly the statement $n\in A^{(k)}$. \hfill $\square$

\subsection*{11.2. A new unconditional upper-density theorem}
For $x\ge 1$, write
\[
A(x):=\#\{n\le x:\ n\in A\},
\qquad
S(x):=\sum_{\substack{n\le x\\ n\in A}}\frac1n.
\]
For each $k\ge 1$, let
\[
A^{(k)}(x):=\#\{n\le x:\ n\in A^{(k)}\},
\qquad
S_k(x):=\sum_{\substack{n\le x\\ n\in A^{(k)}}}\frac1n.
\]

\paragraph{Theorem 11.3 (unconditional upper-density bound by the stage limit).}
For every fixed $k\ge 1$ and all sufficiently large $x$,
\begin{align}
A(x)&\le n_k+A^{(k)}(x),
\tag{11.4}\\
S(x)&\le \sum_{n<n_k}\frac1n+S_k(x).
\tag{11.5}
\end{align}
Consequently,
\begin{equation}
\overline d(A)\le \delta_k,
\qquad
\overline\delta_{\log}(A)\le \delta_k
\qquad (k\ge 1),
\tag{11.6}
\end{equation}
and therefore
\begin{equation}
\overline d(A)\le \delta,
\qquad
\overline\delta_{\log}(A)\le \delta,
\qquad
\delta:=\lim_{k\to\infty}\delta_k.
\tag{11.7}
\end{equation}

\emph{Proof.}
Fix $k\ge 1$.  By Lemma~11.2, if $n\ge n_k$ and $n\in A$, then $n\in A^{(k)}$.  Therefore
\[
A(x)
\le
(n_k-1)+\#\{n_k\le n\le x:\ n\in A^{(k)}\}
\le n_k+A^{(k)}(x),
\]
which is \eqref{11.4}.  Similarly,
\[
S(x)
\le
\sum_{n<n_k}\frac1n
+
\sum_{\substack{n_k\le n\le x\\ n\in A^{(k)}}}\frac1n
\le
\sum_{n<n_k}\frac1n+S_k(x),
\]
which is \eqref{11.5}.

Since $A^{(k)}$ is periodic with density $\delta_k$, periodic asymptotics give
\[
A^{(k)}(x)=\delta_k x+O_k(1),
\qquad
S_k(x)=\delta_k\log x+O_k(1).
\]
Substituting into \eqref{11.4}--\eqref{11.5} and taking $x\to\infty$ yields
\[
\overline d(A)\le \delta_k,
\qquad
\overline\delta_{\log}(A)\le \delta_k.
\]
Since this holds for every $k$ and $\delta_k\downarrow \delta$, we obtain \eqref{11.7}. \hfill $\square$

\paragraph{Corollary 11.4 (corrected Round~8 zero-limit theorem).}
If
\begin{equation}
\delta=\lim_{k\to\infty} d(A^{(k)})=0,
\tag{11.8}
\end{equation}
then $A$ has natural density $0$, and hence logarithmic density $0$.

\emph{Proof.}
By Theorem~11.3,
\[
0\le \overline d(A)\le \delta=0.
\]
Therefore $A$ has natural density $0$.  Since any set of natural density $0$ has logarithmic density $0$ by partial summation, the conclusion follows. \hfill $\square$

\paragraph{Corollary 11.5 (a corrected upper-bound statement for every instance).}
For every instance of Problem~25,
\begin{equation}
\overline\delta_{\log}(A)\le \lim_{k\to\infty} d(A^{(k)}).
\tag{11.9}
\end{equation}
Thus any failure of logarithmic-density existence must come from a \emph{drop below} the stage-density limit, never from an overshoot above it.

\subsection*{11.3. Repair of the Round~7 tail criterion}
The Round~7 theorem was proved from the false global inclusion $A\subseteq A^{(k)}$.  The conclusion remains true, but the comparison must be done only on the tail $[n_k,\infty)$.

\paragraph{Proposition 11.6 (corrected tail criterion).}
For each fixed $k\ge 1$, define
\begin{equation}
\Theta_k:=
\limsup_{x\to\infty}
\frac1{\log x}
\sum_{\substack{i>k\\ n_i\le x}}
\frac{\log(x/n_i)+1}{n_i}.
\tag{11.10}
\end{equation}
If
\begin{equation}
\Theta_k\longrightarrow 0
\qquad (k\to\infty),
\tag{11.11}
\end{equation}
then the logarithmic density $\delta_{\log}(A)$ exists and equals
\begin{equation}
\delta:=\lim_{k\to\infty} d(A^{(k)}).
\tag{11.12}
\end{equation}

\emph{Proof.}
Fix $k\ge 1$ and $x\ge n_k$.  For $n\ge n_k$, if $n\in A^{(k)}\setminus A$, then some later active congruence must fail, so
\[
0\le \mathbf 1_{A^{(k)}}(n)-\mathbf 1_A(n)
\le
\sum_{\substack{i>k\\ n_i\le n}}\mathbf 1_{\,n\equiv a_i\ (\mathrm{mod}\ n_i)}.
\]
Summing from $n_k$ to $x$ with weight $1/n$ gives
\begin{equation}
0
\le
\sum_{\substack{n_k\le n\le x\\ n\in A^{(k)}}}\frac1n
-
\sum_{\substack{n_k\le n\le x\\ n\in A}}\frac1n
\le
\sum_{\substack{i>k\\ n_i\le x}}
\sum_{\substack{n_i\le n\le x\\ n\equiv a_i\ (\mathrm{mod}\ n_i)}}\frac1n.
\tag{11.13}
\end{equation}
The initial segment $[1,n_k)$ contributes only $O_k(1)$ to both $S_k(x)$ and $S(x)$.  Therefore
\begin{equation}
S(x)
\ge
S_k(x)
-
\sum_{\substack{i>k\\ n_i\le x}}
\sum_{\substack{n_i\le n\le x\\ n\equiv a_i\ (\mathrm{mod}\ n_i)}}\frac1n
+O_k(1).
\tag{11.14}
\end{equation}
Using the Round~5 single-progression estimate, we get
\begin{equation}
S(x)
\ge
S_k(x)
-
\sum_{\substack{i>k\\ n_i\le x}}
\left(
\frac1{n_i}\log\frac{x}{n_i}+O\!\left(\frac1{n_i}\right)
\right)
+O_k(1).
\tag{11.15}
\end{equation}
On the other hand, Theorem~11.3 already gives the upper bound
\begin{equation}
S(x)\le S_k(x)+O_k(1).
\tag{11.16}
\end{equation}
Divide \eqref{11.15} and \eqref{11.16} by $\log x$ and let $x\to\infty$.  By definition of $\Theta_k$ and the periodic asymptotic
\[
S_k(x)=\delta_k\log x+O_k(1),
\]
we obtain
\begin{equation}
\delta_k-\Theta_k
\le
\underline\delta_{\log}(A)
\le
\overline\delta_{\log}(A)
\le
\delta_k.
\tag{11.17}
\end{equation}
Now let $k\to\infty$.  Since $\delta_k\to \delta$ and $\Theta_k\to 0$, both the liminf and limsup are squeezed to $\delta$.  Hence $\delta_{\log}(A)=\delta$. \hfill $\square$

\paragraph{Corollary 11.7 (corrected summable-reciprocals theorem).}
If
\begin{equation}
\sum_{i=1}^{\infty}\frac1{n_i}<\infty,
\tag{11.18}
\end{equation}
then $\delta_{\log}(A)$ exists and equals
\begin{equation}
\lim_{k\to\infty} d(A^{(k)}).
\tag{11.19}
\end{equation}

\emph{Proof.}
For each fixed $k$,
\[
\frac1{\log x}
\sum_{\substack{i>k\\ n_i\le x}}
\frac{\log(x/n_i)+1}{n_i}
\le
\sum_{i>k}\frac1{n_i}
+
\frac1{\log x}\sum_{i>k}\frac1{n_i}.
\]
Taking $x\to\infty$ gives
\[
\Theta_k\le \sum_{i>k}\frac1{n_i}\to 0.
\]
Now Proposition~11.6 applies. \hfill $\square$

\subsection*{11.4. A sharpened obstruction: counterexamples can be made tiny-amplitude near $1$}
Let $B$ be any instance with associated lower and upper logarithmic densities
\[
\lambda^-:=\underline\delta_{\log}(B),
\qquad
\lambda^+:=\overline\delta_{\log}(B),
\]
and stage-density limit
\[
\beta:=\lim_{k\to\infty} d(B^{(k)}).
\]
Apply the Round~9 padding transformation with parameter $M\ge 3$, and let $A_M$ denote the padded set.  Write
\[
\alpha_k:=d\bigl(A_M^{(k)}\bigr),
\qquad
\alpha:=\lim_{k\to\infty}\alpha_k.
\]
Round~9/10 give
\begin{align}
\underline\delta_{\log}(A_M)&=\frac{M-2}{M}+\frac{\lambda^-}{M},
\tag{11.20}\\
\overline\delta_{\log}(A_M)&=\frac{M-2}{M}+\frac{\lambda^+}{M},
\tag{11.21}\\
\alpha_k&=\frac{M-2}{M}+\frac{d(B^{(k-1)})}{M}
\qquad (k\ge 1),
\tag{11.22}\\
\alpha&=\frac{M-2}{M}+\frac{\beta}{M}.
\tag{11.23}
\end{align}
Moreover, Theorem~11.3 applied to $A_M$ gives
\begin{equation}
\overline\delta_{\log}(A_M)\le \alpha.
\tag{11.24}
\end{equation}

\paragraph{Theorem 11.8 (small-amplitude high-density obstruction).}
Assume there exists a counterexample to Erd\H{o}s Problem~25.  Then for every $\varepsilon>0$ there exists a counterexample $A$ such that, writing
\[
\delta:=\lim_{k\to\infty} d(A^{(k)}),
\qquad
\lambda^-:=\underline\delta_{\log}(A),
\qquad
\lambda^+:=\overline\delta_{\log}(A),
\]
one has
\begin{align}
1-\varepsilon&<\lambda^-<\lambda^+\le \delta<1,
\tag{11.25}\\
\lambda^+-\lambda^-&<\varepsilon,
\tag{11.26}\\
0\le d(A^{(k)})-\delta&<\varepsilon
\qquad (k\ge 1),
\tag{11.27}\\
\sum_{k\ge 1}\bigl(d(A^{(k)})-d(A^{(k+1)})\bigr)&<\varepsilon.
\tag{11.28}
\end{align}

\emph{Proof.}
Start from any counterexample $B$, so $\lambda^-<\lambda^+$.  Choose an integer $M\ge 3$ so large that
\begin{equation}
\frac{2}{M}<\varepsilon,
\qquad
\frac{1}{M}<\varepsilon.
\tag{11.29}
\end{equation}
Apply the Round~9 padding transformation.  By \eqref{11.20}--\eqref{11.21},
\[
\lambda^-_{A_M}=\frac{M-2+\lambda^-}{M},
\qquad
\lambda^+_{A_M}=\frac{M-2+\lambda^+}{M}.
\]
Since $0\le \lambda^-<\lambda^+\le 1$, we obtain
\[
1-\frac{2}{M}<\lambda^-_{A_M}<\lambda^+_{A_M}<1.
\]
By \eqref{11.29}, this yields the first strict inequalities in \eqref{11.25}.  Also
\[
\lambda^+_{A_M}-\lambda^-_{A_M}=\frac{\lambda^+-\lambda^-}{M}<\frac1M<\varepsilon,
\]
which is \eqref{11.26}.

For the stage densities, Round~10 already gives
\[
0\le \alpha_k-\alpha\le \frac1M<\varepsilon
\qquad (k\ge 1),
\]
which is \eqref{11.27}, and
\[
\sum_{k\ge 1}(\alpha_k-\alpha_{k+1})=\alpha_1-\alpha\le \frac1M<\varepsilon,
\]
which is \eqref{11.28}.

Finally, Theorem~11.3 gives $\lambda^+_{A_M}\le \alpha$, while \eqref{11.23} and $\beta\le 1$ imply
\[
\alpha\le \frac{M-2}{M}+\frac1M=1-\frac1M<1.
\]
Thus
\[
1-\varepsilon<\lambda^-_{A_M}<\lambda^+_{A_M}\le \alpha<1,
\]
which is the remaining part of \eqref{11.25}.  Since logarithmic nonconvergence transfers exactly under padding, $A_M$ is again a counterexample. \hfill $\square$

\paragraph{Interpretation.}
Round~9 showed that the full difficulty survives in an arbitrarily high-density branch, and Round~10 showed it survives when the stage-density sequence is almost constant.  Theorem~11.8 sharpens this substantially:
\begin{quote}
if a counterexample exists, then for every $\varepsilon>0$ there is one whose stage densities, upper logarithmic density, and lower logarithmic density all lie in an interval of width $\varepsilon$ near $1$.
\end{quote}
So the full problem already lives inside an \emph{arbitrarily thin oscillatory window near $1$}.  Any approach that needs a large oscillation gap, noticeable density drop, or substantial total variation in the stage densities cannot settle the full problem on its own.

\section{Adversarial Verification}
We check the new work carefully.
\begin{itemize}
\item \textbf{Concrete flaw identified correctly.}  Lemma~11.1 is an explicit counterexample to the global inclusion $A\subseteq A^{(k)}$.  The issue is exactly the threshold condition $n<n_i$ in the definition of $A$.

\item \textbf{Tail inclusion is the strongest valid comparison.}  Lemma~11.2 uses only that for $n\ge n_k$, all first $k$ constraints are active.  No stronger global inclusion is claimed.

\item \textbf{Upper-density theorem is unconditional.}  Theorem~11.3 uses only the tail inclusion and the periodic asymptotics of each fixed stage set $A^{(k)}$.  No assumptions on growth, compatibility, or reciprocal sums are needed.

\item \textbf{Zero-limit branch repaired correctly.}  Corollary~11.4 follows immediately from the new upper natural-density bound; the earlier flawed proof is no longer needed.

\item \textbf{Corrected Round~7 theorem.}  In Proposition~11.6 the comparison between $A$ and $A^{(k)}$ is performed only on $[n_k,\infty)$, and the initial segment contributes only $O_k(1)$, which disappears after division by $\log x$.

\item \textbf{No circularity in the obstruction theorem.}  Theorem~11.8 uses only previously-vetted Round~9/10 affine padding formulas together with the new unconditional upper bound \eqref{11.24}.  The nonexistence of logarithmic density transfers exactly by affine scaling of the upper/lower densities.

\item \textbf{Compatibility with Round~8 and Round~10.}  There is no contradiction.  Round~8 correctly isolated the benign branch $\delta=0$; the present round shows only that, if a counterexample exists, it can be padded into an almost-stationary branch with all relevant densities near $1$.
\end{itemize}

\section{Final}
\textbf{UNRESOLVED (BUT STRICTLY ADVANCED).}

This round makes two substantive advances.

\begin{enumerate}
\item It identifies and repairs a genuine proof error from Round~7/8: the global inclusion $A\subseteq A^{(k)}$ is false.  Replacing it by the correct tail inclusion yields the new unconditional theorem
\[
\overline d(A)\le \lim_{k\to\infty} d(A^{(k)}),
\qquad
\overline\delta_{\log}(A)\le \lim_{k\to\infty} d(A^{(k)}).
\]
As a corollary, the Round~8 zero-stage-density branch is now proved correctly.

\item It sharpens the obstruction from Round~10.  If a counterexample exists, then for every $\varepsilon>0$ there exists one whose stage densities are uniformly $\varepsilon$-close to a limit $\delta\in(1-\varepsilon,1)$, whose total stage-density variation is $<\varepsilon$, and whose upper and lower logarithmic densities are both in $(1-\varepsilon,1)$ but still distinct, with oscillation gap $<\varepsilon$.
\end{enumerate}

So after Round~11, the unresolved core is even narrower: any counterexample would have to be an \emph{arbitrarily small-amplitude oscillation near density $1$}, while never exceeding the stage-density limit.

\section{Completion Estimate}
\textbf{COMPLETION: 97\%}.

\section{References}
\begin{itemize}
\item No external mathematical references were used in this round.
\item The proofs build directly on the vetted Round~1--Round~10 results supplied in the conversation.
\end{itemize}
